\documentclass[a4paper,11pt]{report}

% =====================================================
% Packages
% =====================================================
\usepackage[utf8]{inputenc}
\usepackage[T1]{fontenc}
\usepackage[french]{babel}
\usepackage{geometry}
\usepackage{listings}
\usepackage{xcolor}
\usepackage{hyperref}

\geometry{margin=2.5cm}

% =====================================================
% Configuration listings C++ & Bash
% =====================================================
\definecolor{codegray}{rgb}{0.95,0.95,0.95}
\definecolor{darkgreen}{rgb}{0.0,0.5,0.0}
\definecolor{darkred}{rgb}{0.6,0.0,0.0}

\lstset{
  language=C++,
  backgroundcolor=\color{codegray},
  basicstyle=\ttfamily\small,
  keywordstyle=\color{blue},
  commentstyle=\color{darkgreen},
  stringstyle=\color{darkred},
  frame=single,
  breaklines=true,
  tabsize=2,
  showstringspaces=false
}

% =====================================================
% Titre
% =====================================================
\title{\textbf{Qualité de code en C++}\\
Bonnes pratiques, robustesse et contexte embarqué}
\author{Document pédagogique complet}

\begin{document}

\maketitle
\tableofcontents
\newpage

% =====================================================
\chapter{Introduction}

\section{Objectif du document}

Ce document a pour objectif de présenter les bonnes pratiques essentielles en langage C++ à travers une approche comparative entre \textbf{mauvais} et \textbf{bon code}. Il vise à améliorer la qualité, la robustesse et la maintenabilité des programmes en mettant en évidence les erreurs courantes et leurs corrections.

\section{Criticité}

Le langage C++ est puissant, expressif et performant, mais il reste exigeant.
Une mauvaise pratique peut rapidement conduire à :
\begin{itemize}
  \item des comportements indéfinis
  \item des erreurs mémoire
  \item des bugs liés aux conversions, aux durées de vie ou au polymorphisme
  \item des défauts difficiles à détecter et à corriger
\end{itemize}

Adopter des règles strictes permet de mieux maîtriser la complexité du langage et de produire un code fiable, lisible et adapté à un usage industriel ou embarqué.

\section{Contenu du document}

Ce document couvre les aspects fondamentaux de la qualité en C++ :
\begin{itemize}
  \item Compilation stricte et gestion des warnings
  \item Bon usage des types, conversions et casts
  \item Gestion sûre des objets, pointeurs, références et durées de vie
  \item Prévention des comportements indéfinis et des erreurs mémoire
  \item Écriture de code lisible, structuré et maintenable
  \item Maîtrise des bonnes pratiques liées à la robustesse, à la concurrence et au contexte embarqué
\end{itemize}

\section{Approche pédagogique}

Chaque chapitre suit une structure simple et efficace :
\begin{itemize}
  \item un exemple de \textbf{mauvais code}
  \item une version corrigée (\textbf{bon code})
  \item une \textbf{explication} claire du problème et de la solution
\end{itemize}

Cette approche permet de comprendre rapidement les erreurs fréquentes et d'adopter les bons réflexes en situation réelle.

% =====================================================
\chapter{Compilation stricte}

\section{Mauvais code}
\begin{lstlisting}
#include <cstdio>

void log_u32(unsigned v) {
    std::printf("%d\n", v);
}
\end{lstlisting}

\section{Bon code}
\begin{lstlisting}
#include <cstdio>

void log_u32(unsigned v) {
    std::printf("%u\n", v);
}
\end{lstlisting}

\section{Explication}
La compilation stricte détecte très tôt les erreurs de format, les conversions implicites dangereuses, les oublis de \texttt{break}, les masquages de variables ou encore les destructeurs non virtuels.
Une base de compilation recommandée est :

\begin{lstlisting}[language=bash]
g++ -std=c++20 -O2 -g -Wall -Wextra -Wpedantic -Werror \
    -Wshadow -Wconversion -Wsign-conversion \
    -Wnull-dereference -Wdouble-promotion -Wformat=2 \
    -Wimplicit-fallthrough -Wold-style-cast -Wnon-virtual-dtor \
    fichier.cpp -o app
\end{lstlisting}

L'objectif est d'obtenir zéro warning.

% =====================================================
\chapter{Allocation dynamique interdite en contexte embarqué}

\section{Mauvais code}
\begin{lstlisting}
int* p = new int(42);
delete p;
\end{lstlisting}

\section{Bon code}
\begin{lstlisting}
#include <array>
#include <cstdint>

static std::array<int, 4> buffer{};
static std::uint32_t value = 42;
\end{lstlisting}

\section{Explication}
En embarqué critique, on évite souvent le heap à cause de la fragmentation, du manque de déterminisme et des échecs d'allocation au runtime.
On préfère des structures statiques ou de taille fixe.

% =====================================================
\chapter{Exceptions souvent interdites}

\section{Mauvais code}
\begin{lstlisting}
int f() {
    throw 42;
}
\end{lstlisting}

\section{Bon code}
\begin{lstlisting}
enum class Status {
    Ok,
    Error
};

Status do_work() noexcept {
    return Status::Ok;
}
\end{lstlisting}

\section{Explication}
Dans beaucoup de projets embarqués, les exceptions sont désactivées pour réduire la taille binaire et garantir un comportement plus déterministe.
On utilise alors des codes retour explicites et \texttt{noexcept}.

% =====================================================
\chapter{Dépassement de tableau}

\section{Mauvais code}
\begin{lstlisting}
int main() {
    int a[3] = {1, 2, 3};
    a[3] = 99;
}
\end{lstlisting}

\section{Bon code}
\begin{lstlisting}
#include <array>

int main() {
    std::array<int, 3> a{1, 2, 3};
    a.at(2) = 99;
}
\end{lstlisting}

\section{Explication}
Accéder hors des bornes d'un tableau provoque un comportement indéfini.
Les conteneurs standards comme \texttt{std::array} permettent des accès plus sûrs, notamment avec \texttt{at()}.

% =====================================================
\chapter{Variables non initialisées}

\section{Mauvais code}
\begin{lstlisting}
#include <cstdio>

int main() {
    int x;
    std::printf("%d\n", x);
}
\end{lstlisting}

\section{Bon code}
\begin{lstlisting}
#include <cstdio>

int main() {
    int x = 0;
    std::printf("%d\n", x);
}
\end{lstlisting}

\section{Explication}
Lire une variable non initialisée donne une valeur indéterminée.
C'est une source fréquente de bugs aléatoires.

% =====================================================
\chapter{Déréférencement de pointeur nul}

\section{Mauvais code}
\begin{lstlisting}
int read(const int* p) {
    return *p;
}
\end{lstlisting}

\section{Bon code}
\begin{lstlisting}
#include <optional>

[[nodiscard]] std::optional<int> read(const int* p) {
    if (!p) {
        return std::nullopt;
    }
    return *p;
}
\end{lstlisting}

\section{Explication}
Déréférencer \texttt{nullptr} provoque un comportement indéfini, souvent un crash.
Il faut vérifier les pointeurs reçus avant usage.

% =====================================================
\chapter{Référence pendante}

\section{Mauvais code}
\begin{lstlisting}
const int& bad_ref() {
    int x = 42;
    return x;
}
\end{lstlisting}

\section{Bon code}
\begin{lstlisting}
int good_value() {
    int x = 42;
    return x;
}
\end{lstlisting}

\section{Explication}
Retourner une référence sur une variable locale est invalide dès la sortie de la fonction.
Il faut retourner une valeur ou utiliser une durée de vie adaptée.

% =====================================================
\chapter{Use-after-scope}

\section{Mauvais code}
\begin{lstlisting}
int* g_ptr = nullptr;

void f() {
    int x = 10;
    g_ptr = &x;
}

int main() {
    f();
    return *g_ptr;
}
\end{lstlisting}

\section{Bon code}
\begin{lstlisting}
void set_value(int& out) {
    out = 10;
}

int main() {
    int x = 0;
    set_value(x);
    return x;
}
\end{lstlisting}

\section{Explication}
Conserver l'adresse d'une variable locale au-delà de sa portée mène à un comportement indéfini.
Il faut laisser l'appelant gérer la durée de vie des objets qu'il fournit.

% =====================================================
\chapter{Overflow signé}

\section{Mauvais code}
\begin{lstlisting}
#include <climits>

int main() {
    int x = INT_MAX;
    x += 1;
}
\end{lstlisting}

\section{Bon code}
\begin{lstlisting}
#include <limits>
#include <optional>

[[nodiscard]] std::optional<int> add_safe(int a, int b) {
    if (b > 0 && a > std::numeric_limits<int>::max() - b) {
        return std::nullopt;
    }
    if (b < 0 && a < std::numeric_limits<int>::min() - b) {
        return std::nullopt;
    }
    return a + b;
}
\end{lstlisting}

\section{Explication}
Le dépassement de capacité d'un entier signé est un comportement indéfini en C++.
Il faut contrôler les bornes avant l'opération.

% =====================================================
\chapter{Division par zéro}

\section{Mauvais code}
\begin{lstlisting}
int main() {
    int a = 10;
    int b = 0;
    return a / b;
}
\end{lstlisting}

\section{Bon code}
\begin{lstlisting}
#include <optional>

[[nodiscard]] std::optional<int> div_checked(int a, int b) {
    if (b == 0) {
        return std::nullopt;
    }
    return a / b;
}
\end{lstlisting}

\section{Explication}
Diviser un entier par zéro est un comportement indéfini.
Le contrôle préalable est obligatoire.

% =====================================================
\chapter{Conversions signé / non signé}

\section{Mauvais code}
\begin{lstlisting}
#include <cstddef>

int main() {
    int i = -1;
    std::size_t n = static_cast<std::size_t>(i);
    (void)n;
}
\end{lstlisting}

\section{Bon code}
\begin{lstlisting}
#include <optional>
#include <cstddef>

[[nodiscard]] std::optional<std::size_t> to_size(int i) {
    if (i < 0) {
        return std::nullopt;
    }
    return static_cast<std::size_t>(i);
}
\end{lstlisting}

\section{Explication}
Un entier négatif converti en type non signé devient une très grande valeur.
Il faut vérifier le domaine de validité avant conversion.

% =====================================================
\chapter{Narrowing et troncature}

\section{Mauvais code}
\begin{lstlisting}
#include <cstdint>

std::uint8_t v = 300;
\end{lstlisting}

\section{Bon code}
\begin{lstlisting}
#include <optional>
#include <cstdint>

[[nodiscard]] std::optional<std::uint8_t> to_u8(int value) {
    if (value < 0 || value > 255) {
        return std::nullopt;
    }
    return static_cast<std::uint8_t>(value);
}
\end{lstlisting}

\section{Explication}
Affecter une valeur trop grande dans un type plus petit tronque l'information.
Cette troncature doit être évitée ou validée explicitement.

% =====================================================
\chapter{Double promotion}

\section{Mauvais code}
\begin{lstlisting}
float x = 1.0f;
float y = x + 0.1;
\end{lstlisting}

\section{Bon code}
\begin{lstlisting}
float x = 1.0f;
float y = x + 0.1f;
\end{lstlisting}

\section{Explication}
Un littéral décimal comme \texttt{0.1} est un \texttt{double} par défaut.
Cela force une promotion implicite qui peut dégrader les performances ou provoquer des warnings.

% =====================================================
\chapter{Cast C-style dangereux}

\section{Mauvais code}
\begin{lstlisting}
double d = 3.14;
int x = (int&)d;
\end{lstlisting}

\section{Bon code}
\begin{lstlisting}
double d = 3.14;
int x = static_cast<int>(d);
\end{lstlisting}

\section{Explication}
Les casts C-style masquent plusieurs types de conversion possibles.
Les casts C++ rendent l'intention explicite et permettent une lecture plus sûre.

% =====================================================
\chapter{Affectation au lieu de comparaison}

\section{Mauvais code}
\begin{lstlisting}
bool ok = false;

if (ok = true) {
}
\end{lstlisting}

\section{Bon code}
\begin{lstlisting}
bool ok = false;

if (ok) {
}
\end{lstlisting}

\section{Explication}
Utiliser \texttt{=} dans une condition modifie la variable.
Cette erreur classique peut être détectée avec les warnings de compilation.

% =====================================================
\chapter{Switch avec fallthrough involontaire}

\section{Mauvais code}
\begin{lstlisting}
int f(int x) {
    switch (x) {
        case 0:
            doA();
        case 1:
            doB();
            break;
        default:
            return 0;
    }
    return 1;
}
\end{lstlisting}

\section{Bon code}
\begin{lstlisting}
int f(int x) {
    switch (x) {
        case 0:
            doA();
            break;
        case 1:
            doB();
            break;
        default:
            return 0;
    }
    return 1;
}
\end{lstlisting}

\section{Explication}
Sans \texttt{break}, l'exécution continue dans le cas suivant.
Cela peut être voulu, mais dans ce cas il faut le documenter explicitement.

% =====================================================
\chapter{Switch non exhaustif sur enum}

\section{Mauvais code}
\begin{lstlisting}
enum class State {
    Idle,
    Run,
    Error
};

void handle(State s) {
    switch (s) {
        case State::Idle:
            break;
        case State::Run:
            break;
    }
}
\end{lstlisting}

\section{Bon code}
\begin{lstlisting}
enum class State {
    Idle,
    Run,
    Error
};

void handle(State s) {
    switch (s) {
        case State::Idle:
            break;
        case State::Run:
            break;
        case State::Error:
            break;
    }
}
\end{lstlisting}

\section{Explication}
Oublier un état dans un \texttt{switch} sur une énumération crée un bug silencieux.
Chaque valeur doit être traitée explicitement ou via un \texttt{default} justifié.

% =====================================================
\chapter{Accolades manquantes}

\section{Mauvais code}
\begin{lstlisting}
if (ok)
    doA();
    doB();
\end{lstlisting}

\section{Bon code}
\begin{lstlisting}
if (ok) {
    doA();
    doB();
}
\end{lstlisting}

\section{Explication}
Sans accolades, l'indentation peut devenir trompeuse.
Il faut préférer des blocs explicites, même pour un seul statement.

% =====================================================
\chapter{Comparaison flottante directe}

\section{Mauvais code}
\begin{lstlisting}
float a = 0.1f + 0.2f;

if (a == 0.3f) {
}
\end{lstlisting}

\section{Bon code}
\begin{lstlisting}
#include <cmath>

bool nearly_equal(float a, float b, float eps) {
    return std::fabs(a - b) <= eps;
}
\end{lstlisting}

\section{Explication}
Les flottants sont approximatifs.
Une comparaison directe avec \texttt{==} est donc fragile dans la plupart des cas.

% =====================================================
\chapter{Commentaires trompeurs}

\section{Mauvais code}
\begin{lstlisting}
#include <array>
#include <cstdint>

// Allocate 128 bytes of RAM
std::array<std::uint8_t, 14> txBuffer;
\end{lstlisting}

\section{Bon code}
\begin{lstlisting}
#include <array>
#include <cstdint>

// txBuffer contient le message a emettre (14 octets fixes).
std::array<std::uint8_t, 14> txBuffer{};
\end{lstlisting}

\section{Explication}
Un commentaire faux ou imprécis nuit à la compréhension.
Un bon commentaire explique la contrainte, l'intention ou la raison métier.

% =====================================================
\chapter{Convention de nommage}

\section{Mauvais code}
\begin{lstlisting}
int MaxValue = 100;
int max_value = 200;
int MAXVALUE = 300;

void DoStuff();
void do_stuff();
\end{lstlisting}

\section{Bon code}
\begin{lstlisting}
constexpr int kMaxValue = 100;
int current_value = 0;

struct PacketHeader {
    int id;
};

void process_packet(const PacketHeader& h);
\end{lstlisting}

\section{Explication}
Une convention cohérente améliore fortement la lisibilité.
Il faut définir une règle simple et la respecter dans tout le projet.

% =====================================================
\chapter{Fonctions trop longues}

\section{Mauvais code}
\begin{lstlisting}
void handle_request(Request r) {
    /* parse + validate + auth + log + business + format */
}
\end{lstlisting}

\section{Bon code}
\begin{lstlisting}
Response handle_request(const Request& r) {
    const auto parsed = parse(r);

    if (!validate(parsed)) {
        return error("bad request");
    }

    const auto user = authenticate(parsed);
    return execute_business(user, parsed);
}
\end{lstlisting}

\section{Explication}
Une fonction qui fait trop de choses devient difficile à tester, relire et maintenir.
Il faut découper les responsabilités.

% =====================================================
\chapter{Nombres magiques}

\section{Mauvais code}
\begin{lstlisting}
int timeout_ms = 1500;
int retries = 4;

if (retries > 3) {
    abort();
}
\end{lstlisting}

\section{Bon code}
\begin{lstlisting}
constexpr int kTimeoutMs = 1500;
constexpr int kMaxRetries = 3;

int retries = 4;
int timeout_ms = kTimeoutMs;

if (retries > kMaxRetries) {
    abort();
}
\end{lstlisting}

\section{Explication}
Les valeurs magiques rendent le code opaque.
Les constantes nommées rendent l'intention explicite.

% =====================================================
\chapter{Valeur de retour ignorée}

\section{Mauvais code}
\begin{lstlisting}
struct Frame {
    int id;
};

bool send_frame(const Frame& f);

void task(const Frame& f) {
    send_frame(f);
}
\end{lstlisting}

\section{Bon code}
\begin{lstlisting}
struct Frame {
    int id;
};

[[nodiscard]] bool send_frame(const Frame& f);

void task(const Frame& f) {
    if (!send_frame(f)) {
        /* gerer l'erreur */
    }
}
\end{lstlisting}

\section{Explication}
Ignorer une valeur de retour masque une erreur potentielle.
\texttt{[[nodiscard]]} aide à imposer cette discipline.

% =====================================================
\chapter{Const-correctness}

\section{Mauvais code}
\begin{lstlisting}
#include <array>

int sum(std::array<int,3>& a) {
    return a[0] + a[1] + a[2];
}
\end{lstlisting}

\section{Bon code}
\begin{lstlisting}
#include <array>

int sum(const std::array<int,3>& a) {
    return a[0] + a[1] + a[2];
}
\end{lstlisting}

\section{Explication}
Une API qui n'a pas besoin de modifier son argument doit l'exprimer avec \texttt{const}.
Cela clarifie le contrat et élargit les cas d'usage.

% =====================================================
\chapter{Concurrence : data race}

\section{Mauvais code}
\begin{lstlisting}
#include <thread>

int x = 0;

void f() {
    for (int i = 0; i < 100000; i++) {
        x++;
    }
}
\end{lstlisting}

\section{Bon code}
\begin{lstlisting}
#include <thread>
#include <atomic>

std::atomic<int> x{0};

void f() {
    for (int i = 0; i < 100000; i++) {
        x.fetch_add(1);
    }
}
\end{lstlisting}

\section{Explication}
Deux threads qui modifient la même variable sans synchronisation créent une course de données.
Le résultat devient non déterministe.

% =====================================================
\chapter{Destructeur non virtuel}

\section{Mauvais code}
\begin{lstlisting}
struct Base {
    ~Base() = default;
};

struct Der : Base {
    ~Der() = default;
};
\end{lstlisting}

\section{Bon code}
\begin{lstlisting}
struct Base {
    virtual ~Base() = default;
};

struct Der : Base {
    ~Der() override = default;
};
\end{lstlisting}

\section{Explication}
Si une classe est utilisée polymorphiquement, son destructeur doit être virtuel.
Sinon la destruction via pointeur de base est dangereuse.

% =====================================================
\chapter{Object slicing}

\section{Mauvais code}
\begin{lstlisting}
#include <cstdio>

struct Base {
    virtual ~Base() = default;
    virtual void f() { std::printf("B\n"); }
};

struct Der : Base {
    void f() override { std::printf("D\n"); }
    int x = 1;
};

int main() {
    Der d;
    Base b = d;
    b.f();
}
\end{lstlisting}

\section{Bon code}
\begin{lstlisting}
#include <cstdio>

struct Base {
    virtual ~Base() = default;
    virtual void f() { std::printf("B\n"); }
};

struct Der : Base {
    void f() override { std::printf("D\n"); }
    int x = 1;
};

int main() {
    Der d;
    Base& b = d;
    b.f();
}
\end{lstlisting}

\section{Explication}
Copier un objet dérivé dans un objet base par valeur supprime la partie dérivée.
Il faut passer par référence ou pointeur quand on veut conserver le polymorphisme.

% =====================================================
\chapter{Buffer overflow via sprintf}

\section{Mauvais code}
\begin{lstlisting}
#include <cstdio>

int main() {
    char buf[8];
    std::sprintf(buf, "value=%d", 123456);
}
\end{lstlisting}

\section{Bon code}
\begin{lstlisting}
#include <cstdio>
#include <cstddef>

bool make(char* out, std::size_t cap, int v) {
    int n = std::snprintf(out, cap, "value=%d", v);
    return n >= 0 && static_cast<std::size_t>(n) < cap;
}
\end{lstlisting}

\section{Explication}
\texttt{sprintf} n'impose aucune limite de taille.
\texttt{snprintf} permet d'éviter les dépassements de tampon.

% =====================================================
\chapter{Variables d'interruption et volatile}

\section{Mauvais code}
\begin{lstlisting}
bool data_ready = false;

void UART_ISR() {
    data_ready = true;
}

int main() {
    while (!data_ready) {
    }
}
\end{lstlisting}

\section{Bon code}
\begin{lstlisting}
volatile bool data_ready = false;

void UART_ISR() {
    data_ready = true;
}

int main() {
    while (!data_ready) {
    }
}
\end{lstlisting}

\section{Explication}
Une variable modifiée par interruption peut être optimisée à tort par le compilateur.
Dans ce cas, \texttt{volatile} indique qu'il faut relire la valeur en mémoire.

% =====================================================
\chapter{Initialisation globale risquée}

\section{Mauvais code}
\begin{lstlisting}
/* A.cpp */
HardwareTimer timerA;

/* B.cpp */
TimerManager manager(timerA);
\end{lstlisting}

\section{Bon code}
\begin{lstlisting}
constinit int g_hardware_config = 42;

HardwareTimer& get_timer() {
    static HardwareTimer timer;
    return timer;
}
\end{lstlisting}

\section{Explication}
L'ordre d'initialisation des objets globaux répartis dans plusieurs unités de traduction n'est pas maîtrisé.
Le pattern \emph{construct on first use} ou \texttt{constinit} réduit ce risque.

% =====================================================
\chapter{Alignement mémoire et padding}

\section{Mauvais code}
\begin{lstlisting}
#include <cstdint>

struct HardwareReg {
    std::uint8_t flag;
    std::uint32_t data;
};
\end{lstlisting}

\section{Bon code}
\begin{lstlisting}
#include <cstdint>

struct __attribute__((packed)) HardwareRegPacked {
    std::uint8_t flag;
    std::uint32_t data;
};

static_assert(sizeof(HardwareRegPacked) == 5);
\end{lstlisting}

\section{Explication}
Le compilateur peut insérer des octets de bourrage dans une structure.
Quand on mappe une structure sur du matériel ou un protocole, il faut maîtriser sa taille exacte.

% =====================================================
\chapter{Conclusion}

Un code C++ de qualité professionnelle en approche DevSecOps repose sur :
\begin{itemize}
    \item une compilation stricte sans warning ;
    \item une gestion rigoureuse des conversions, de la mémoire et des durées de vie ;
    \item des API explicites avec \texttt{const}, \texttt{[[nodiscard]]} et \texttt{noexcept} ;
    \item une attention particulière au polymorphisme, à la concurrence et à la robustesse ;
    \item un style cohérent et une structuration claire des responsabilités ;
    \item une adaptation aux contraintes embarquées quand le contexte l'exige.
\end{itemize}

Ces pratiques permettent de produire un code plus fiable, plus lisible, plus maintenable et plus sûr dans la durée.

\end{document}